\chapter{Concept Generation}
% this is some introduction about the topic
\section{Pre-class Activities}
\subsection{Activity 1: Create Concept Classification Trees OR Concept Combination Tables}
\textbf{Objective:} \\

The objective of this activity is to walk you in developing either concept classification trees and concept combination tables. These tools help to better organize and categorize the identified solution concepts for each critical sub-problem, and also help to explore different combinations of sub-solutions.

\textbf{Procedure:}
\begin{enumerate}
    \item Start this activity by first familiarizing yourself regarding the purpose and benefits of concept classification trees and concept combination tables.
    \item For each critical sub-problem identified in activity 1, develop concept classification trees
    \begin{enumerate}
        \item Label the nodes or branches to indicate the different categories or attributes used for classification. Example of categories or attributes including functionality, cost, feasibility, user experience
        \item Assign each solution concept to the appropriate category or attribute within the concept classification tree.
    \end{enumerate}
    \item Develop concept combination tables to systematically explore different combinations of sub-solutions.
    \begin{enumerate}
        \item Create a table with
        \begin{itemize}
            \item columns representing the sub-problems
            \item rows representing the different solution concepts.
        \end{itemize}
        \item Identify and develop possible combinations by cross-referencing the sub-problems and solution concepts.
    \end{enumerate}
\end{enumerate}

\subsection{Activity 2: Generate rough product concepts } % if available
\textbf{Objective:} \\
The objective of this activity is to generate rough product concepts by combining the solution concepts from different sub-problems within the concept combination table

\textbf{Procedure:}
\begin{enumerate}
    \item Refer to the output from activity 1, and generate rough product concepts by combining the solution concepts from different sub-problems within the concept combination table
    \begin{enumerate}
        \item Be creative and explore various combinations to generate a range of potential product concepts.
        \item Provide brief description to effectively convey the key aspects of the proposed rough product concepts.
        \item Provide information such as the technology utilized, working principles, intended functionality, and unique features of the concept.
    \end{enumerate}
    \item Repeat steps 1 for at least five generated concepts that show promise or align well with the project goals for further development.
\end{enumerate}

\section{In-class Activities}
\subsubsection{Activity 1: Create Concept combination table}
Compile all findings and discuss with all members whether it is more advantageous to select the best concept combination table that has been drafted by one of the members, or to combine elements from multiple concept combination tables to create a new concept combination table that systematically explores different combinations of sub-solutions. Once consensus is reached, draft a concept combination table that reflects the collective agreement of all team members.

\subsubsection{Activity 2: Generate rough product concepts} % if available
Compile all findings and discuss with all members about the generated concepts that have been drafted by each one of the members. Select the top 5 generated concepts that show promise or align well with the project goals for further development.